\backhead{FURTHER EXERCISES}{khred}
\begin{exlist}
\item A lateral of 80 m, inside diameter 16 mm, carries 80 emitters of 2 L\,h$^{-1}$. Find the friction loss with the Christiansen factor ($C=140$).\exans{0.16 m}
\item An emitter has $q=kH^{x}$ with $q=2.0$ L\,h$^{-1}$ at 10 m and $x=0.5$. Find $q$ at 8 m and 12 m and the variation of the flow over this range of head.\exans{1.79 and 2.19 L\,h$^{-1}$; 20 per cent}
\item A drip block has 2400 emitters of 2 L\,h$^{-1}$. Find the flow of the block in m$^3$\,h$^{-1}$ and in L\,s$^{-1}$. If the crop needs 4 mm\,d$^{-1}$ on 0.6 ha, how long must the block run per day at an efficiency of 0.9?\exans{4.8 m$^3$\,h$^{-1}$; 1.33 L\,s$^{-1}$; 5.6 h}
\item Find $EU$ for emitters with $C_v=0.07$, two per plant, and $q_{\min}/q_{\mathrm{avg}}=0.9$.\exans{84.3 per cent}
\item A wetted bulb in a sandy loam has a radius of 0.35 m at the surface. What emitter spacing gives a continuous strip with 20 per cent overlap?\exans{0.56 m}
\end{exlist}
\backhead{ESSAY QUESTIONS}{khred}
{\small\itshape Answer each question in about one page, with equations, a sketch and at least one number from this chapter.}\par\vspace{3pt}
\begin{exlist}
\item Explain the hydraulics of a drip lateral and how the length, the diameter and the slope determine the uniformity.
\item Discuss the wetted bulb: what controls its shape and why it matters for the spacing of emitters and for the crop.
\item Compare drip and sprinkler irrigation in terms of evaporation, drainage and salinity.
\item Explain emission uniformity and how manufacturing variation and pressure variation combine.
\item Discuss the design of a drip system for a whole farm: blocks, mains, filtration and the pump.
\item Why does drip irrigation not always save water at the scale of a basin?
\end{exlist}
\dykband{EXERCISES IN SWAFL}
{\small Open SWAFL (desktop or online) with the benchmark project. The chapter of the \href{https://github.com/halilunaziru73-creator/Coupled-Water-Fluxes-Pore-Scale-Dynamics-Precision-Irrigation-Mulching-Feedback-Loops/raw/main/Software/v1.1.0/SWAFL_Desktop/docs/manual_v1.1.0/SWAFL_Manual_v1.1.0.pdf}{SWAFL manual} that explains each task is given in brackets; record your values in a table and compare them with what is expected (\emph{Check}).}\par\vspace{3pt}
\begin{exlist}
\item {\color{khblue}[Manual, Chapter 12]} Open \emph{Drip irrigation design}; for the loam and the sand, record the bulb radius and depth for the same emitter.\par\noindent{\small\itshape Check: the sand bulb is narrower and deeper.}
\item {\color{khblue}[Manual, Chapter 12]} Change the lateral length in steps of 20 m and record the emission uniformity. Find the longest acceptable lateral.\par\noindent{\small\itshape Check: the uniformity falls with length.}
\item {\color{khblue}[Manual, Chapter 12]} Use \emph{To size a pump} and record the head and the power for the farm.\par\noindent{\small\itshape Check: the pump head includes friction and pressure.}
\item {\color{khblue}[Manual, Chapter 18]} Compare the three systems on this field and record which one the program prefers and why.\par\noindent{\small\itshape Check: the comparison lists water, energy and uniformity.}
\end{exlist}
